% GENERATED FILE. Edit canonical lessons, then run npm run generate:books.
% canonical-source-hash: fnv1a64:ffb425ed488ff2d5
% canonical-lessons: ML-C497-tiradesam, ML-C497-alimukham, ML-C497-veliyerram, ML-C497-veliyirakkam, ML-C497-katalakramanam

\chapter{Along the Coast}
\label{ch:ml-along-the-coast}
\hlchaptermodality{497}

\begin{chapteropening}
\textbf{By the end of this chapter:} \emph{I can say the coastal region, an estuary, high tide, low tide and a sea surge.}

Complete the last lesson of chapter 497: I can say the coastal region, an estuary, high tide, low tide and a sea surge.
\end{chapteropening}

\section[\texorpdfstring{\ml{തീരദേശം} (tīradēśaṁ)}{തീരദേശം (tīradēśaṁ)}]{\ml{തീരദേശം} (tīradēśaṁ) — the coastal region}

\label{lesson:ML-C497-tiradesam}



\begin{quote}
Before the new one: say the Malayalam for a dam, then the Malayalam for to dry up.
\end{quote}

\subsection*{You'll want to know: \ml{തീരദേശം}}
\textbf{\ml{തീരദേശം}} — \emph{tīradēśaṁ} — \textquotedblleft{}the coastal region\textquotedblright{}.

\begin{quote}
\ml{തീരദേശത്ത്} \ml{ധാരാളം} \ml{മീൻപിടുത്തക്കാർ} \ml{താമസിക്കുന്നു}.
\end{quote}

\begin{quote}
\emph{tīradēśattŭ dhārāḷaṁ mīnpiṭuttakkār tāmasikkunnu.}
\end{quote}

\begin{quote}
Many fishing people live in the coastal region.
\end{quote}

\ml{തീരം}, a shore, plus \ml{ദേശം}, a land.

\begin{cousinweb}
\textbf{\ml{തീരദേശം}} (\emph{tīradēśaṁ}) is a Malayalam compound built from \ml{തീരം}, a shore, from Sanskrit \emph{tīra} (CDIAL 5842), and \ml{ദേശം}, a land, from Sanskrit \emph{deśa} (Monier-Williams).
\end{cousinweb}

\begin{grammarlens}[title={What you've built}]
One more word for getting around.
\end{grammarlens}

\subsection*{Your turn}
\begin{itemize}
\raggedright
  \item \emph{Say it:} \emph{tīradēśaṁ}
  \item \emph{Say it:} \emph{tīradēśaṁ}, once more
  \item \emph{Say it:} \emph{tīradēśattŭ dhārāḷaṁ mīnpiṭuttakkār tāmasikkunnu.} — Many fishing people live in the coastal region.
  \item \emph{Recall:} say the Malayalam for a dam, then the Malayalam for to dry up, then say \emph{tīradēśaṁ} again
\end{itemize}

\begin{tcolorbox}[breakable,colback=teal!4,colframe=teal!35!black,title={Before you move on}]
What does \ml{തീരദേശം} mean? (\textquotedblleft{}The coastal region\textquotedblright{}.) Say it once more.
\end{tcolorbox}
\section[\texorpdfstring{\ml{അഴിമുഖം} (aḻimukhaṁ)}{അഴിമുഖം (aḻimukhaṁ)}]{\ml{അഴിമുഖം} (aḻimukhaṁ) — an estuary; a river mouth}

\label{lesson:ML-C497-alimukham}



\begin{quote}
Before the new one: say the Malayalam for to dry up, then the Malayalam for the coastal region.
\end{quote}

\subsection*{You'll want to know: \ml{അഴിമുഖം}}
\textbf{\ml{അഴിമുഖം}} — \emph{aḻimukhaṁ} — \textquotedblleft{}an estuary; a river mouth\textquotedblright{}.

\begin{quote}
\ml{പുഴ} \ml{കടലിൽ} \ml{ചേരുന്നത്} \ml{അഴിമുഖത്താണ്}.
\end{quote}

\begin{quote}
\emph{puḻa kaṭalil cērunnatŭ aḻimukhattāṇŭ.}
\end{quote}

\begin{quote}
The river meets the sea at the estuary.
\end{quote}

\ml{അഴി}, an opening to the sea, plus \ml{മുഖം}, a face: the face of the river towards the sea.

\begin{cousinweb}
\textbf{\ml{അഴിമുഖം}} (\emph{aḻimukhaṁ}) is a Malayalam compound built from \ml{അഴി}, a bar or opening where a river meets the sea, whose origin is not settled, and \ml{മുഖം}, a face, from Sanskrit \emph{mukha} (CDIAL 10158).
\end{cousinweb}

\begin{grammarlens}[title={What you've built}]
One more word for getting around.
\end{grammarlens}

\subsection*{Your turn}
\begin{itemize}
\raggedright
  \item \emph{Say it:} \emph{aḻimukhaṁ}
  \item \emph{Say it:} \emph{aḻimukhaṁ}, once more
  \item \emph{Say it:} \emph{puḻa kaṭalil cērunnatŭ aḻimukhattāṇŭ.} — The river meets the sea at the estuary.
  \item \emph{Recall:} say the Malayalam for to dry up, then the Malayalam for the coastal region, then say \emph{aḻimukhaṁ} again
\end{itemize}

\begin{tcolorbox}[breakable,colback=teal!4,colframe=teal!35!black,title={Before you move on}]
What does \ml{അഴിമുഖം} mean? (\textquotedblleft{}An estuary; a river mouth\textquotedblright{}.) Say it once more.
\end{tcolorbox}
\section[\texorpdfstring{\ml{വേലിയേറ്റം} (vēliyēṟṟaṁ)}{വേലിയേറ്റം (vēliyēṟṟaṁ)}]{\ml{വേലിയേറ്റം} (vēliyēṟṟaṁ) — high tide}

\label{lesson:ML-C497-veliyerram}



\begin{quote}
Before the new one: say the Malayalam for the coastal region, then the Malayalam for an estuary.
\end{quote}

\subsection*{You'll want to know: \ml{വേലിയേറ്റം}}
\textbf{\ml{വേലിയേറ്റം}} — \emph{vēliyēṟṟaṁ} — \textquotedblleft{}high tide\textquotedblright{}.

\begin{quote}
\ml{വേലിയേറ്റത്തിൽ} \ml{കടൽവെള്ളം} \ml{കായലിലേക്ക്} \ml{കയറി}.
\end{quote}

\begin{quote}
\emph{vēliyēṟṟattil kaṭalveḷḷaṁ kāyalilēkkŭ kayaṟi.}
\end{quote}

\begin{quote}
At high tide sea water came into the backwaters.
\end{quote}

\ml{വേലി}, the tide, plus \ml{ഏറ്റം}, a rise: the sea comes far up the shore.

\begin{cousinweb}
\textbf{\ml{വേലിയേറ്റം}} (\emph{vēliyēṟṟaṁ}) is a Malayalam compound built from \ml{വേലി}, the tide, usually traced to Sanskrit \emph{velā}, \textquotedblleft{}the boundary of sea and land\textquotedblright{} (Monier-Williams), and the native \ml{ഏറ്റം}, a rise.
\end{cousinweb}

\begin{grammarlens}[title={What you've built}]
One more word for getting around.
\end{grammarlens}

\subsection*{Your turn}
\begin{itemize}
\raggedright
  \item \emph{Say it:} \emph{vēliyēṟṟaṁ}
  \item \emph{Say it:} \emph{vēliyēṟṟaṁ}, once more
  \item \emph{Say it:} \emph{vēliyēṟṟattil kaṭalveḷḷaṁ kāyalilēkkŭ kayaṟi.} — At high tide sea water came into the backwaters.
  \item \emph{Recall:} say the Malayalam for the coastal region, then the Malayalam for an estuary, then say \emph{vēliyēṟṟaṁ} again
\end{itemize}

\begin{tcolorbox}[breakable,colback=teal!4,colframe=teal!35!black,title={Before you move on}]
What does \ml{വേലിയേറ്റം} mean? (\textquotedblleft{}High tide\textquotedblright{}.) Say it once more.
\end{tcolorbox}
\section[\texorpdfstring{\ml{വേലിയിറക്കം} (vēliyiṟakkaṁ)}{വേലിയിറക്കം (vēliyiṟakkaṁ)}]{\ml{വേലിയിറക്കം} (vēliyiṟakkaṁ) — low tide}

\label{lesson:ML-C497-veliyirakkam}



\begin{quote}
Before the new one: say the Malayalam for an estuary, then the Malayalam for high tide.
\end{quote}

\subsection*{You'll want to know: \ml{വേലിയിറക്കം}}
\textbf{\ml{വേലിയിറക്കം}} — \emph{vēliyiṟakkaṁ} — \textquotedblleft{}low tide\textquotedblright{}.

\begin{quote}
\ml{വേലിയിറക്കത്തിൽ} \ml{കുട്ടികൾ} \ml{മണലിൽ} \ml{ഞണ്ടുകളെ} \ml{പിടിച്ചു}.
\end{quote}

\begin{quote}
\emph{vēliyiṟakkattil kuṭṭikaḷ maṇalil ñaṇṭukaḷe piṭiccu.}
\end{quote}

\begin{quote}
At low tide the children caught crabs on the sand.
\end{quote}

\ml{ഇറക്കം}, a going down, is the noun of \ml{ഇറങ്ങുക}.

\begin{cousinweb}
\textbf{\ml{വേലിയിറക്കം}} (\emph{vēliyiṟakkaṁ}) is a Malayalam compound built from \ml{വേലി}, the tide, usually traced to Sanskrit \emph{velā}, and the native \ml{ഇറക്കം}, a going down (DEDR 516).
\end{cousinweb}

\begin{grammarlens}[title={What you've built}]
One more word for getting around.
\end{grammarlens}

\subsection*{Your turn}
\begin{itemize}
\raggedright
  \item \emph{Say it:} \emph{vēliyiṟakkaṁ}
  \item \emph{Say it:} \emph{vēliyiṟakkaṁ}, once more
  \item \emph{Say it:} \emph{vēliyiṟakkattil kuṭṭikaḷ maṇalil ñaṇṭukaḷe piṭiccu.} — At low tide the children caught crabs on the sand.
  \item \emph{Recall:} say the Malayalam for an estuary, then the Malayalam for high tide, then say \emph{vēliyiṟakkaṁ} again
\end{itemize}

\begin{tcolorbox}[breakable,colback=teal!4,colframe=teal!35!black,title={Before you move on}]
What does \ml{വേലിയിറക്കം} mean? (\textquotedblleft{}Low tide\textquotedblright{}.) Say it once more.
\end{tcolorbox}
\section[\texorpdfstring{\ml{കടലാക്രമണം} (kaṭalākramaṇaṁ)}{കടലാക്രമണം (kaṭalākramaṇaṁ)}]{\ml{കടലാക്രമണം} (kaṭalākramaṇaṁ) — a sea surge; coastal erosion}

\label{lesson:ML-C497-katalakramanam}



\begin{quote}
Before the new one: say the Malayalam for high tide, then the Malayalam for low tide.
\end{quote}

\subsection*{You'll want to know: \ml{കടലാക്രമണം}}
\textbf{\ml{കടലാക്രമണം}} — \emph{kaṭalākramaṇaṁ} — \textquotedblleft{}a sea surge; coastal erosion\textquotedblright{}.

\begin{quote}
\ml{കടലാക്രമണത്തിൽ} \ml{തീരത്തെ} \ml{വീടുകൾ} \ml{തകർന്നു}.
\end{quote}

\begin{quote}
\emph{kaṭalākramaṇattil tīratte vīṭukaḷ takarnnu.}
\end{quote}

\begin{quote}
In the sea surge the houses on the shore were destroyed.
\end{quote}

Literally an attack by the sea: \ml{കടൽ} plus \ml{ആക്രമണം}, an attack. A word of the monsoon news.

\begin{cousinweb}
\textbf{\ml{കടലാക്രമണം}} (\emph{kaṭalākramaṇaṁ}) is a Malayalam compound built from the native \ml{കടൽ}, the sea (DEDR 1118), and \ml{ആക്രമണം}, an attack, from Sanskrit \emph{ākramaṇa}, \textquotedblleft{}stepping upon, attacking\textquotedblright{} (Monier-Williams).
\end{cousinweb}

\begin{grammarlens}[title={What you've built}]
One more word for getting around.
\end{grammarlens}

\subsection*{Your turn}
\begin{itemize}
\raggedright
  \item \emph{Say it:} \emph{kaṭalākramaṇaṁ}
  \item \emph{Say it:} \emph{kaṭalākramaṇaṁ}, once more
  \item \emph{Say it:} \emph{kaṭalākramaṇattil tīratte vīṭukaḷ takarnnu.} — In the sea surge the houses on the shore were destroyed.
  \item \emph{Recall:} say the Malayalam for high tide, then the Malayalam for low tide, then say \emph{kaṭalākramaṇaṁ} again
\end{itemize}

\begin{tcolorbox}[breakable,colback=teal!4,colframe=teal!35!black,title={Before you move on}]
What does \ml{കടലാക്രമണം} mean? (\textquotedblleft{}A sea surge; coastal erosion\textquotedblright{}.) Say it once more.
\end{tcolorbox}
